\section{Mengen}
\label{sec:fundamentals_sets}

\begin{DefinitionBox}{Menge}{fundamentals_set}
  Eine \emph{Menge} \EnglishTerm{set} ist eine Zusammenfassung bestimmter,
  aber unterschiedlicher, \enquote{Elemente} zu einem Ganzen.
\end{DefinitionBox}

\begin{RemarkBox}
  Ein Element kann alles sein, was wir uns vorstellen können, zum Beispiel
  eine Zahl, ein Objekt oder auch eine Menge. Jedes Element kann höchstens
  einmal in einer Menge enthalten sein.
\end{RemarkBox}

Für Mengen werden zwei Schreibweisen verwendet:
\begin{enumerate}[label=(\roman*),itemsep=0.8em]
  \item \emph{Aufzählende Form:} Eine Menge, die die Elemente $a$, $b$ und
    $c$ enthält, schreibt man als
    \begin{ClozeArea}[\linewidth]{1.5cm}
      \[
        M=\{a,b,c\}.
      \]
    \end{ClozeArea}
    Die Reihenfolge spielt dabei keine Rolle.

  \item \emph{Beschreibende Form:} Eine Menge, die alle Elemente mit einer
    bestimmten Eigenschaft enthält, schreibt man als
    \begin{ClozeArea}[\linewidth]{1.5cm}
      \[
        M=\{x\mid x\text{ hat eine bestimmte Eigenschaft}\}.
      \]
    \end{ClozeArea}
\end{enumerate}

\begin{ExampleBox}{Aufzählende vs. Beschreibende Form}{fundamentals_set_notation_example}
  Die Menge $M$ der ganzen Zahlen zwischen $1$ und $5$ kann dauf 2 Arbeiten geschrieben werden:
      \[
        \Cloze{8cm}{\{n \in \mathbb{N} \mid 1 \leq n \leq 5 \} = M = \{1,2,3,4,5\}}.
      \]
\end{ExampleBox}


\begin{DefinitionBox}{Leere Menge}{fundamentals_emptyset}
  Unter der \emph{leere Menge} \EnglishTerm{empty set} versteht man eine Menge, die kein Element enthält. Sie wird mit dem Symbol $\emptyset$ bezeichnet.
\end{DefinitionBox}




\begin{DefinitionBox}{Element einer Menge}{fundamentals_set_element}
  Ist $a$ ein \emph{Element} \EnglishTerm{element} aus der Menge $M$, so
  schreibt man
  \begin{align*}
    &a\in M
    \qquad \text{\enquote{ist Element von}}, \\
    %
    \text{sonst} \quad
    &a\notin M
    \qquad\text{\enquote{ist kein Element von}}.
  \end{align*}
\end{DefinitionBox}



\begin{ExampleBox}{Elemente von Mengen}{fundamentals_elements_of_sets}
  \begin{enumerate}[label=(\alph*),itemsep=0.6em]
    \item $1\ \Cloze{1.5cm}{\in}\ \{1,2,3\}$
    \item $\text{Peter}\ \Cloze{1.5cm}{\notin}\
      \{\text{Franz},\text{Anne}\}$
  \end{enumerate}
\end{ExampleBox}

\subsection{Teilmengen}
\label{subsec:fundamentals_subsets}

\begin{DefinitionBox}{Teilmenge}{fundamentals_subset}
  Eine Menge $A$ heißt \emph{Teilmenge} \EnglishTerm{subset} einer Menge $B$,
  wenn jedes Element von $A$ auch Element von $B$ ist. Man schreibt
  \[
    A\subseteq B
    \quad\Longleftrightarrow\quad
    \forall x:\; x\in A\Rightarrow x\in B.
  \]
\end{DefinitionBox}

\begin{RemarkBox}[title=\Lang{Bemerkung}{Remark}]
  \begin{minipage}[t]{0.46\linewidth}
    \textbf{Teilmenge}
    \[
      B\subseteq A
    \]
    Hier ist $B=A$ erlaubt.
  \end{minipage}
  \hfill
  \begin{minipage}[t]{0.46\linewidth}
    \textbf{Echte Teilmenge}
    \[
      B\subsetneq A
    \]
    Hier ist $B=A$ nicht erlaubt, es gilt also $B\neq A$.
  \end{minipage}
\end{RemarkBox}

Zwei Mengen sind genau dann gleich, wenn sie dieselben Elemente enthalten.
Damit gilt
\[
  A=B
  \quad\Longleftrightarrow\quad
  A\subseteq B\ \land\ B\subseteq A.
\]

\begin{ExampleBox}{Teilmengen}{fundamentals_subset_example}
  Für $A=\{1,2\}$ und $B=\{1,2,3\}$ gilt
  \[
    A\ \Cloze{1.8cm}{\subseteq B},
    \qquad
    B\ \Cloze{1.8cm}{\not\subseteq A}.
  \]
\end{ExampleBox}

\clearpage
\subsection{Mengenoperationen}
\label{subsec:fundamentals_set_operations}

Zwischen Mengen existieren Operationen, mit denen aus gegebenen Mengen neue
Mengen gebildet werden können. Zu den wichtigsten gehören Schnitt, Vereinigung
und Differenz.

\subsubsection{Schnittmenge}
\label{subsubsec:fundamentals_set_intersection}

\noindent\begin{minipage}{\linewidth}
\begin{DefinitionBox}{Schnittmenge}{fundamentals_set_intersection}
  \begin{minipage}[c]{0.70\linewidth}
    Die \emph{Schnittmenge} \EnglishTerm{intersection} zweier Mengen $A$ und
    $B$ enthält alle Elemente, die sowohl in der Menge $A$ als auch in der
    Menge $B$ enthalten sind.
    \begin{ClozeArea}[24]{3}
      \[
        A\cap B=\{x\mid x\in A\land x\in B\}.
      \]
    \end{ClozeArea}
  \end{minipage}\hfill
  \begin{minipage}[c]{0.25\linewidth}
    \centering
    \begin{tikzpicture}[scale=1.2]
  \begin{scope}
    \clip (-0.5,0) circle (0.75);
    \fill[THMBlue!30] (0.5,0) circle (0.75);
  \end{scope}
  \draw[thick] (-0.5,0) circle (0.75);
  \draw[thick] (0.5,0) circle (0.75);
  \node at (-0.5,0) {$A$};
  \node at (0.5,0) {$B$};
\end{tikzpicture}

  \end{minipage}
\end{DefinitionBox}

\begin{ExampleBox}{Schnittmenge}{fundamentals_set_intersection_example}
  \begin{enumerate}[label=(\alph*),itemsep=0.6em]
    \item $\{1,2\}\cap\{2,3\}=\Cloze{3cm}{\{2\}}$
    \item $\{1,2,3\}\cap\{1,2,3\}=\Cloze{3cm}{\{1,2,3\}}$
    \item $M\cap\emptyset=\Cloze{2.5cm}{\emptyset}\quad(\forall M)$
  \end{enumerate}
\end{ExampleBox}
\end{minipage}



\subsubsection{Vereinigungsmenge}
\label{subsubsec:fundamentals_set_union}

\noindent\begin{minipage}{\linewidth}
\begin{DefinitionBox}{Vereinigungsmenge}{fundamentals_set_union}
  \begin{minipage}[c]{0.70\linewidth}
    Die \emph{Vereinigungsmenge} \EnglishTerm{union} zweier Mengen $A$ und
    $B$ enthält alle Elemente, die in der Menge $A$, in der Menge $B$ oder in
    beiden Mengen enthalten sind.
    \begin{ClozeArea}[24]{3}
      \[
        A\cup B=\{x\mid x\in A\lor x\in B\}.
      \]
    \end{ClozeArea}
  \end{minipage}\hfill
  \begin{minipage}[c]{0.25\linewidth}
    \centering
    \begin{tikzpicture}[scale=1.2]
  \fill[THMBlue!30] (-0.5,0) circle (0.75);
  \fill[THMBlue!30] (0.5,0) circle (0.75);
  \draw[thick] (-0.5,0) circle (0.75);
  \draw[thick] (0.5,0) circle (0.75);
  \node at (-0.5,0) {$A$};
  \node at (0.5,0) {$B$};
\end{tikzpicture}

  \end{minipage}
\end{DefinitionBox}

\begin{ExampleBox}{Vereinigungsmenge}{fundamentals_set_union_example}
  \begin{enumerate}[label=(\alph*),itemsep=0.6em]
    \item $\{1,2\}\cup\{2,3\}=\Cloze{3cm}{\{1,2,3\}}$
    \item $\{1,2,3\}\cup\{1,2,3\}=\Cloze{3cm}{\{1,2,3\}}$
    \item $M\cup\emptyset=\Cloze{2.5cm}{M}\quad(\forall M)$
  \end{enumerate}
\end{ExampleBox}
\end{minipage}


\subsubsection{Differenzmenge}
\label{subsubsec:fundamentals_set_difference}

\noindent\begin{minipage}{\linewidth}
\begin{DefinitionBox}{Differenzmenge}{fundamentals_set_difference}
  \begin{minipage}[c]{0.70\linewidth}
    Die \emph{Differenzmenge} \EnglishTerm{set difference} $A\setminus B$
    enthält alle Elemente, die in der Menge $A$, aber nicht in der Menge $B$
    enthalten sind.
    \begin{ClozeArea}[24]{3}
      \[
        A\setminus B=\{x\mid x\in A\land x\notin B\}.
      \]
    \end{ClozeArea}
  \end{minipage}\hfill
  \begin{minipage}[c]{0.25\linewidth}
    \centering
    \begin{tikzpicture}[scale=1.2]
  \begin{scope}
    \clip (-0.5,0) circle (0.75);
    \fill[THMBlue!30] (-0.5,0) circle (0.75);
    \clip (0.5,0) circle (0.75);
    \fill[white] (0.5,0) circle (0.75);
  \end{scope}
  \draw[thick] (-0.5,0) circle (0.75);
  \draw[thick] (0.5,0) circle (0.75);
  \node at (-0.5,0) {$A$};
  \node at (0.5,0) {$B$};
\end{tikzpicture}

  \end{minipage}
\end{DefinitionBox}

\begin{ExampleBox}{Differenzmenge}{fundamentals_set_difference_example}
  \begin{enumerate}[label=(\alph*),itemsep=0.6em]
    \item $\{1,2\}\setminus\{2,3\}=\Cloze{3cm}{\{1\}}$
    \item $\{1,2,3\}\setminus\{1,2,3\}=\Cloze{2.5cm}{\emptyset}$
    \item $M \setminus \emptyset = \Cloze{2.5cm}{M}$
  \end{enumerate}
\end{ExampleBox}
\end{minipage}

